\chapter{2005年全国I卷（理）}
\section{单选题}
\begin{problem}
复数 \(\frac{\sqrt{2} - \i^3}{1 - \sqrt{2}\i} =\)（\quad）

\choices
    {\(\i\)}
    {\(-\i\)}
    {\(2\sqrt{2} - \i\)}
    {\(-2\sqrt{2} + \i\)}
\answer{A}

\begin{solution}
因为 \(\i^3=-\i\)，所以
\(\sqrt{2}-\i^3=\sqrt{2}+\i\). 分子分母同乘分母的共轭复数 \(1+\sqrt{2}\i\)，得
\(\displaystyle\frac{\sqrt{2}+\i}{1-\sqrt{2}\i}=\frac{(\sqrt{2}+\i)(1+\sqrt{2}\i)}{1+2}=\frac{3\i}{3}=\i\). 故选 A.
\end{solution}
\end{problem}

\begin{problem}
设 \(I\) 为全集，\(S_{1}\)，\(S_{2}\)，\(S_{3}\) 是 \(I\) 的三个非空子集且 \(S_{1} \cup S_{2} \cup S_{3} = I\)，则下面论断正确的是（\quad）

\choices
    {\(\complement_{I}S_{1}\cap (S_{2} \cup S_{3}) = \varnothing\)}
    {\(S_{1}\subseteq (\complement_{I}S_{2}\cap \complement_{I}S_{3})\)}
    {\(\complement_{I}S_{1}\cap \complement_{I}S_{2}\cap \complement_{I}S_{3} = \varnothing\)}
    {\(S_{1}\subseteq (\complement_{I}S_{2}\cup \complement_{I}S_{3})\)}
\answer{C}

\begin{solution}
由 \(S_{1}\cup S_{2}\cup S_{3}=I\)，取补集并利用德 Morgan 定律，得
\(\complement_{I}S_{1}\cap\complement_{I}S_{2}\cap\complement_{I}S_{3}=\complement_{I}(S_{1}\cup S_{2}\cup S_{3})=\complement_{I}I=\varnothing\). 因此正确的论断是 \(\circled{3}\)，故选 C.
\end{solution}
\end{problem}

\begin{problem}
一个与球心距离为 \(1\) 的平面截球所得的圆面面积为 \(\pi\)，则球的表面积为（\quad）

\choices
    {\(8\sqrt{2}\pi\)}
    {\(8\pi\)}
    {\(4\sqrt{2}\pi\)}
    {\(4\pi\)}
\answer{B}

\begin{solution}
设球的半径为 \(R\)，截得圆面的半径为 \(r\). 由圆面面积为 \(\pi\)，得 \(r=1\). 设球心为 \(O\)，截面圆心为 \(H\)，并在截面圆周上取一点 \(R_0\)，则 \(OH=1\)，且 \(R_0H\perp OH\)，所以在直角三角形 \(OHR_0\) 中，\(R^{2}=OH^{2}+R_0H^{2}=1^{2}+1^{2}=2\). 球的表面积为 \(4\pi R^{2}=8\pi\)，故选 B.
\end{solution}
\end{problem}

\begin{problem}
已知直线 \(l\) 过点 \((-2,0)\)，当直线 \(l\) 与圆 \(y=1-\sqrt{1-x^{2}}\)（\(-1\leqslant x\leqslant1\)）有两个交点时，其斜率 \(k\) 的取值范围是（\quad）

\choices
    {\((-2\sqrt{2}, 2\sqrt{2})\)}
    {\((-\sqrt{2}, \sqrt{2})\)}
    {\(\left(\frac{\sqrt{2}}{4}, \frac{\sqrt{2}}{4}\right)\)}
    {\(\left(-\frac{1}{8}, \frac{1}{8}\right)\)}
\answer{\(0<k\leqslant\frac{1}{3}\)（题面与选项不一致）}

\begin{solution}
直线 \(l\) 的方程为 \(y=k(x+2)\). 题面所给曲线的函数值满足 \(0\leqslant y\leqslant1\)，而 \(-1\leqslant x\leqslant1\) 时 \(x+2>0\)，所以若有两个交点，必有 \(k>0\). 令
\(g(x)=1-\sqrt{1-x^{2}}-k(x+2)\)，\(-1\leqslant x\leqslant1\).
曲线与直线的交点正好是 \(g(x)=0\) 的根. 在 \(-1<x<1\) 内，
\[
g''(x)=\frac{1}{(1-x^{2})^{3/2}}>0
\]
所以 \(g\) 严格凸，因而至多有两个零点. 又
\(g(-1)=1-k\)，\(g(0)=-2k\)，\(g(1)=1-3k\).
当 \(0<k<\frac13\) 时，\(g(-1)>0\)、\(g(0)<0\)、\(g(1)>0\)，由介值定理在 \((-1,0)\) 和 \((0,1)\) 各有一个零点，恰有两个交点. 当 \(k=\frac13\) 时，\(g(-1)>0\)、\(g(0)<0\)、\(g(1)=0\)，交点为一个内点和端点 \((1,1)\)，仍有两个不同交点. 当 \(\frac13<k<1\) 时，\(g(0)<0\)、\(g(1)<0\)，严格凸性说明 \((0,1)\) 内无零点，左侧至多一个交点；当 \(k\geqslant1\) 时，\(g(-1)\leqslant0\)、\(g(1)<0\)，也不可能有两个交点. \(k=0\) 或 \(k<0\) 同样不满足两个交点条件.

因此，按原试卷打印的曲线，斜率范围为 \(0<k\leqslant\frac13\)，与题中四个选项均不一致；该题的原试卷题面、选项与公开标准版本存在冲突，故按题面不能对应合法选项.
\end{solution}
\end{problem}

\begin{problem}
如图，在多面体 \(ABCDEF\) 中，已知 \(ABCD\) 是边长为 \(1\) 的正方形，且 \(\triangle ADE\)，\(\triangle BCF\) 均为正三角形，\(EF \parallel AB\)，\(EF = 2\)，则该多面体的体积为（\quad）


\bitmapfigure[max width=0.48\linewidth]{img/2005/national_paper_1_science/q5_fig1.png}

\choices
    {\(\frac{\sqrt{2}}{3}\)}
    {\(\frac{\sqrt{3}}{3}\)}
    {\(\frac{4}{3}\)}
    {\(\frac{3}{2}\)}
\answer{A}

\begin{solution}
建立空间直角坐标系，取正方形 \(ABCD\) 所在平面为 \(z=0\)，令
\(A(0,0,0)\)，\(B(1,0,0)\)，\(C(1,1,0)\)，\(D(0,1,0)\). 设
\(E=(u,\frac12,h)\)，则由 \(\triangle ADE\) 为正三角形，得
\(u^{2}+\frac14+h^{2}=1\). 由于 \(EF\parallel AB\) 且 \(EF=2\)，可设 \(F=(u+2,\frac12,h)\). 由 \(\triangle BCF\) 为正三角形，得
\((u+1)^{2}+\frac14+h^{2}=1\). 两式相减得 \(u=-\frac12\)，再代回得 \(h^{2}=\frac12\). 取 \(h=\frac{1}{\sqrt{2}}\).

记 \(z=th\)（\(0\leqslant t\leqslant1\)）的水平截面为 \(K_t\). 该截面是底面正方形与顶面线段 \(EF\) 按比例 \(1-t:t\) 的凸组合，因此其两边长分别为
\((1-t)\cdot1+t\cdot2=1+t\) 和 \((1-t)\cdot1=1-t\)，从而
\(S_{K_t}=1-t^{2}\). 所以多面体体积为
\(V=\displaystyle\int_{0}^{h}\left(1-\frac{z^{2}}{h^{2}}\right)\,\mathrm dz=\frac{2h}{3}=\frac{\sqrt{2}}{3}\). 故选 A.
\end{solution}
\end{problem}

\begin{problem}
已知双曲线 \(\frac{x^2}{a^2} - y^2 = 1\)（\(a > 0\)）的一条准线与抛物线 \(y^2 = -6x\) 的准线重合，则该双曲线的离心率为（\quad）

\choices
    {\(\frac{\sqrt{3}}{2}\)}
    {\(\frac{3}{2}\)}
    {\(\frac{\sqrt{6}}{2}\)}
    {\(\frac{2\sqrt{3}}{3}\)}
\answer{D}

\begin{solution}
双曲线 \(\displaystyle\frac{x^{2}}{a^{2}}-y^{2}=1\) 的半焦距满足 \(c^{2}=a^{2}+1\)，离心率为 \(e=\frac{c}{a}\)，其准线方程为 \(x=\pm\frac{a}{e}=\pm\frac{a^{2}}{c}\). 抛物线 \(y^{2}=-6x\) 写成 \(y^{2}=2px\)，其中 \(p=-3\)，故其准线为 \(x=-\frac p2=\frac32\). 因而
\(\displaystyle\frac{a^{2}}{c}=\frac32\). 令 \(t=a^{2}\)，平方得 \(\displaystyle\frac{t^{2}}{t+1}=\frac94\)，即 \(4t^{2}-9t-9=0\). 由于 \(t>0\)，得 \(t=3\)，于是 \(c=2\)，\(e=\frac{2}{\sqrt3}=\frac{2\sqrt3}{3}\). 故选 D.
\end{solution}
\end{problem}

\begin{problem}
当 \(0 < x < \frac{\pi}{2}\) 时，函数 \(f(x) = \frac{1 + \cos 2x + 8\sin^2 x}{\sin 2x}\) 的最小值为（\quad）

\choices
    {\(2\)}
    {\(2\sqrt{3}\)}
    {\(4\)}
    {\(4\sqrt{3}\)}
\answer{C}

\begin{solution}
当 \(0<x<\frac\pi2\) 时，\(\tan x>0\)，且
\(\displaystyle f(x)=\frac{2\cos^{2}x+8\sin^{2}x}{2\sin x\cos x}=\frac{1+4\tan^{2}x}{\tan x}=\frac1{\tan x}+4\tan x\).
令 \(t=\tan x>0\)，由基本不等式
\(\displaystyle\frac1t+4t\geqslant2\sqrt{\frac1t\cdot4t}=4\)，等号当且仅当 \(t=\frac12\) 时成立. 因此最小值为 \(4\)，故选 C.
\end{solution}
\end{problem}

\begin{problem}
设 \(b > 0\)，二次函数 \(y = ax^{2} + bx + a^{2} - 1\) 的图象为下列之一，则 \(a\) 的值为（\quad）

\choices
    {\choicebitmap[width=0.15\paperwidth]{img/2005/national_paper_1_science/q8_optA.png}}
    {\choicebitmap[width=0.15\paperwidth]{img/2005/national_paper_1_science/q8_optB.png}}
    {\choicebitmap[width=0.15\paperwidth]{img/2005/national_paper_1_science/q8_optC.png}}
    {\choicebitmap[width=0.15\paperwidth]{img/2005/national_paper_1_science/q8_optD.png}}
\answer{C}

\begin{solution}
抛物线与 \(y\) 轴的交点纵坐标为 \(a^{2}-1\)，开口方向由 \(a\) 的符号决定，对称轴为 \(x=-\frac{b}{2a}\). 题图中前两幅图的对称轴为 \(y\) 轴，与 \(b>0\) 矛盾；第四幅图经过原点且在原点处向下穿过 \(x\) 轴，因而其一次项系数应为负，也与 \(b>0\) 矛盾.

第三幅图开口向下且经过原点，所以 \(a^{2}-1=0\) 且 \(a<0\)，从而 \(a=-1\). 这时图象在原点处的斜率为 \(b>0\)，与第三幅图相符，故选 C.
\end{solution}
\end{problem}

\begin{problem}
设 \(0 < a < 1\)，函数 \(f(x) = \log_a(a^{2x} - 2a^x - 2)\)，则使 \(f(x) < 0\) 的 \(x\) 取值范围是（\quad）

\choices
    {\(f(x)\)}
    {\((0, +\infty)\)}
    {\(\left(-\infty, \log_a 3\right)\)}
    {\((\log_a 3, +\infty)\)}
\answer{C}

\begin{solution}
令 \(t=a^{x}>0\). 由于 \(0<a<1\)，函数 \(\log_{a}u\) 的定义域要求 \(u>0\)，且 \(\log_{a}u<0\) 等价于 \(u>1\). 因而
\(f(x)<0\iff a^{2x}-2a^{x}-2>1\iff t^{2}-2t-3>0\iff (t-3)(t+1)>0\).
因为 \(t>0\)，所以 \(t>3\)，即 \(a^{x}>3\). 由 \(0<a<1\) 时指数函数单调递减，得 \(x<\log_{a}3\). 这个条件自动保证真数 \(t^{2}-2t-2>0\)，故取值范围为 \((-\infty,\log_{a}3)\)，选 C.
\end{solution}
\end{problem}

\begin{problem}
在坐标平面上，不等式组 \(\begin{cases}
y\geqslant x-1,\\
y\leqslant-3|x|+1
\end{cases}\) 所表示的平面区域的面积为（\quad）

\choices
    {\(\sqrt{2}\)}
    {\(\frac{3}{2}\)}
    {\(\frac{3\sqrt{2}}{2}\)}
    {\(2\)}
\answer{B}

\begin{solution}
当 \(x\leqslant0\) 时，上界为 \(y=3x+1\)，下界为 \(y=x-1\)，两直线交于 \((-1,-2)\)，并在 \(x=0\) 处分别取 \(1\) 和 \(-1\). 这部分区域是底为 \(1\)、两端高为 \(0\) 和 \(2\) 的梯形，面积为 \(\frac{0+2}{2}\cdot1=1\).

当 \(x\geqslant0\) 时，上界为 \(y=-3x+1\)，与下界 \(y=x-1\) 交于 \((\frac12,-\frac12)\). 这部分区域与左侧合并后应直接按分段积分计算：
\(\displaystyle S=\int_{-1}^{0}\bigl[(3x+1)-(x-1)\bigr]\,\mathrm dx+\int_{0}^{1/2}\bigl[(-3x+1)-(x-1)\bigr]\,\mathrm dx=1+\frac12=\frac32\).
因此平面区域面积为 \(\frac32\)，故选 B.
\end{solution}
\end{problem}

\begin{problem}
在 \(\triangle ABC\) 中，已知 \(\tan \frac{A + B}{2} = \sin C\)，给出以下四个论断：

\begin{circlelist}
\item \(\tan A \cdot \cot B = 1\)；
\item \(0 < \sin A + \sin B \leqslant \sqrt{2}\)；
\item \(\sin^2 A + \cos^2 B = 1\)；
\item \(\cos^2 A + \cos^2 B = \sin^2 C\).
\end{circlelist}

其中正确的是（\quad）

\choices
    {\(\circled{1}\circled{3}\)}
    {\(\circled{2}\circled{4}\)}
    {\(\circled{1}\circled{4}\)}
    {\(\circled{2}\circled{3}\)}
\answer{B}

\begin{solution}
在三角形中 \(A+B=\pi-C\)，所以
\(\displaystyle\tan\frac{A+B}{2}=\cot\frac C2=\sin C=2\sin\frac C2\cos\frac C2\). 因为 \(0<C<\pi\)，故 \(\sin\frac C2\) 与 \(\cos\frac C2\) 均为正，化简得 \(2\sin^{2}\frac C2=1\)，从而 \(C=\frac\pi2\)，即 \(A+B=\frac\pi2\).

于是 \(\sin A+\sin B=2\sin\frac{A+B}{2}\cos\frac{A-B}{2}=\sqrt2\cos\frac{A-B}{2}\). 由于 \(A\)，\(B>0\)，有 \(0<\cos\frac{A-B}{2}\leqslant1\)，故 \(\circled{2}\) 正确. 又 \(\cos B=\sin A\)，所以 \(\cos^{2}A+\cos^{2}B=\cos^{2}A+\sin^{2}A=1=\sin^{2}C\)，故 \(\circled{4}\) 正确.

例如取 \(A=\frac\pi6\)，\(B=\frac\pi3\)，则 \(C=\frac\pi2\)，题设条件成立，但 \(\tan A\cdot\cot B=\frac13\ne1\)，且 \(\sin^{2}A+\cos^{2}B=\frac12\ne1\)，故 \(\circled{1}\)、\(\circled{3}\) 不正确. 因此选 B.
\end{solution}
\end{problem}

\begin{problem}
过三棱柱任意两个顶点的直线共\(15\)条，其中异面直线有（\quad）

\choices
    {\(18\text{ 对}\)}
    {\(24\text{ 对}\)}
    {\(30\text{ 对}\)}
    {\(36\text{ 对}\)}
\answer{D}

\begin{solution}
三棱柱的顶点记为 \(A\)，\(B\)，\(C\)，\(A'\)，\(B'\)，\(C'\). 任取两个顶点确定的直线共有 \(\mathrm C_{6}^{2}=15\) 条，可分为三条侧棱、六条三角形底面的边和六条侧面的对角线三类.

以侧棱 \(AA'\) 为例，与它异面的直线为 \(BC\)、\(BC'\)、\(CB'\)、\(B'C'\)，共 \(4\) 条. 以底面边 \(AB\) 为例，与它异面的直线为 \(CA'\)、\(CB'\)、\(CC'\)、\(A'C'\)、\(B'C'\)，共 \(5\) 条. 以侧面上的对角线 \(AB'\) 为例，与它异面的直线为 \(BC\)、\(BC'\)、\(CA'\)、\(CC'\)、\(A'C'\)，共 \(5\) 条. 其余同类直线由三棱柱的对称性得到相同数量. 因而把每一对异面直线计算两次，异面直线对数为
\(\displaystyle\frac{3\times4+6\times5+6\times5}{2}=36\). 故选 D.
\end{solution}
\end{problem}

\section{填空题}
\begin{problem}
若正整数 \(m\) 满足 \(10^{m-1} < 2^{512} < 10^{m}\)，则 \(m =\fillinblank{}\).（\(\lg 2 \approx 0.3010\)）
\answer{155}

\begin{solution}
对不等式取常用对数，得 \(m-1<512\lg2<m\). 由 \(\lg2\approx0.3010\)，有 \(512\lg2\approx154.112\)，所以 \(m-1<154.112<m\). 因为 \(m\) 为正整数，故 \(m=155\).
\end{solution}
\end{problem}

\begin{problem}
\(\left(2x - \frac{1}{\sqrt{x}}\right)^9\) 的展开式中，常数项为\fillinblank{}.（用数字作答）
\answer{672}

\begin{solution}
展开式的通项为 \(\displaystyle\mathrm C_{9}^{k}(2x)^{9-k}\left(-\frac1{\sqrt{x}}\right)^{k}\)，其中 \(k=0\)，\(1\)，\(\ldots\)，\(9\). 其 \(x\) 的指数为 \(9-k-\frac{k}{2}=9-\frac{3k}{2}\). 要得到常数项，须有 \(9-\frac{3k}{2}=0\)，故 \(k=6\). 常数项为
\(\displaystyle\mathrm C_{9}^{6}2^{3}(-1)^{6}=84\times8=672\).
\end{solution}
\end{problem}

\begin{problem}
\(\triangle ABC\) 的外接圆的圆心为 \(O\)，两条边上的高的交点为 \(H\)，\(\overrightarrow{OH} = m(\overrightarrow{OA} + \overrightarrow{OB} + \overrightarrow{OC})\)，则实数 \(m =\)\fillinblank{}.
\answer{1}

\begin{solution}
以 \(O\) 为原点，记 \(\overrightarrow{OA}=\bs a\)，\(\overrightarrow{OB}=\bs b\)，\(\overrightarrow{OC}=\bs c\). 因为 \(A\)，\(B\)，\(C\) 在以 \(O\) 为圆心的圆上，所以 \(|\bs a|=|\bs b|=|\bs c|\). 令 \(\bs h=\bs a+\bs b+\bs c\)，取点 \(H_0\) 使 \(\overrightarrow{OH_0}=\bs h\). 则
\(\overrightarrow{AH_0}=\bs h-\bs a=\bs b+\bs c\)，\(\overrightarrow{BC}=\bs c-\bs b\)，从而 \(\overrightarrow{AH_0}\cdot\overrightarrow{BC}=|\bs c|^{2}-|\bs b|^{2}=0\). 同理可得 \(BH_0\perp AC\)，所以点 \(H_0\) 是三角形的垂心，即 \(H_0=H\)，从而
\(\overrightarrow{OH}=\bs a+\bs b+\bs c=\overrightarrow{OA}+\overrightarrow{OB}+\overrightarrow{OC}\).
因此 \((m-1)\bs h=\bs0\). 当 \(\bs h\ne\bs0\)（即 \(ABC\) 不是等边三角形）时，\(m=1\)，故按题目预期答案填 1. 若 \(ABC\) 为等边三角形，则 \(\bs h=\bs0\)，题设等式对任意 \(m\) 都成立；题面未排除这一特殊情形.
\end{solution}
\end{problem}

\begin{problem}
在正方体 \(ABCD - A'B'C'D'\) 中，过对角线 \(BD'\) 的一个平面交 \(AA'\) 于 \(E\)，交 \(CC'\) 于 \(F\)，则：

\begin{circlelist}
\item 四边形 \(BFD'E\) 一定是平行四边形；
\item 四边形 \(BFD'E\) 有可能是正方形；
\item 四边形 \(BFD'E\) 在底面 \(ABCD\) 内的投影一定是正方形；
\item 平面 \(BFD'E\) 有可能垂直于平面 \(BB'D\).
\end{circlelist}

以上结论正确的为\fillinblank{}.（写出所有正确结论的编号）
\answer{\(\circled{1}\circled{3}\circled{4}\)}

\begin{solution}
设正方体棱长为 \(a\)，取坐标
\(A(0,0,0)\)，\(B(a,0,0)\)，\(C(a,a,0)\)，\(D(0,a,0)\)，\(A'(0,0,a)\)，\(D'(0,a,a)\). 设 \(E=(0,0,ta)\)，其中 \(0\leqslant t\leqslant1\). 过 \(B\)，\(D'\)，\(E\) 的平面方程可写为
\(tx+ty-y+z-ta=0\). 令 \(x=y=a\)，得到 \(F=(a,a,(1-t)a)\).

于是 \(\overrightarrow{BF}=(0,a,(1-t)a)=\overrightarrow{ED'}\)，且 \(\overrightarrow{FD'}=(-a,0,ta)=\overrightarrow{BE}\)，所以四边形 \(BFD'E\) 一定是平行四边形，\(\circled{1}\) 正确. 它在底面 \(ABCD\) 上的投影四个顶点为 \(B\)，\(C\)，\(D\)，\(A\)，故投影一定是正方形，\(\circled{3}\) 正确.

相邻两边的内积为 \(\overrightarrow{BF}\cdot\overrightarrow{FD'}=a^{2}t(1-t)\). 若平行四边形为正方形，须有 \(t=0\) 或 \(t=1\)，但此时两邻边长度分别为 \(a\) 与 \(a\sqrt2\)，不相等，所以 \(\circled{2}\) 不正确.

平面 \(BFD'E\) 的一个法向量可取 \((t,t-1,1)\)，平面 \(BB'D\) 的一个法向量可取 \((1,1,0)\). 当 \(t=\frac12\) 时两法向量的内积为 \(0\)，故此时两平面垂直，\(\circled{4}\) 正确. 所以正确结论为 \(\circled{1}\)、\(\circled{3}\)、\(\circled{4}\).
\end{solution}
\end{problem}

\section{解答题}
\begin{problem}
设函数 \(f(x) = \sin(2x + \varphi)\)（\(-\pi < \varphi < 0\)），\(y = f(x)\) 图象的一条对称轴是直线 \(x = \frac{\pi}{8}\).

\begin{enumerate}
\item 求 \(\varphi\)；
\item 求函数 \(y = f(x)\) 的单调增区间；
\item 证明直线 \(5x - 2y + c = 0\) 与函数 \(y = f(x)\) 的图象不相切.
\end{enumerate}
\end{problem}

\begin{solution}
\begin{enumerate}
\item 函数 \(y=\sin(2x+\varphi)\) 的对称轴满足 \(2x+\varphi=\dfrac{\pi}{2}+k\pi\)，其中 \(k\in\Z\). 代入 \(x=\dfrac{\pi}{8}\)，得 \(\dfrac{\pi}{4}+\varphi=\dfrac{\pi}{2}+k\pi\)，所以 \(\varphi=\dfrac{\pi}{4}+k\pi\). 由 \(-\pi<\varphi<0\) 得 \(k=-1\)，故 \(\varphi=-\dfrac{3\pi}{4}\).
\item 此时 \(f(x)=\sin\left(2x-\dfrac{3\pi}{4}\right)\)，且 \(f'(x)=2\cos\left(2x-\dfrac{3\pi}{4}\right)\). 当 \(f'(x)\geqslant0\) 时，函数单调增加，即 \(-\dfrac{\pi}{2}+2k\pi\leqslant2x-\dfrac{3\pi}{4}\leqslant\dfrac{\pi}{2}+2k\pi\)，所以单调增区间为 \(\left[k\pi+\dfrac{\pi}{8},\,k\pi+\dfrac{5\pi}{8}\right]\)，其中 \(k\in\Z\).
\item 若直线 \(5x-2y+c=0\) 与函数图象相切于点 \((x_0,f(x_0))\)，则直线斜率应等于函数在该点的导数，即 \(f'(x_0)=\dfrac{5}{2}\). 但对任意 \(x_0\)，都有 \(\left|f'(x_0)\right|=\left|2\cos\left(2x_0-\dfrac{3\pi}{4}\right)\right|\leqslant2<\dfrac{5}{2}\)，矛盾. 因此直线 \(5x-2y+c=0\) 与函数图象不相切.
\end{enumerate}
\end{solution}

\begin{problem}
已知四棱锥 \(P - ABCD\) 的底面为直角梯形，\(AB \parallel DC\)，\(\angle DAB = 90^{\circ}\)，\(PA \perp\) 底面 \(ABCD\)，且 \(PA = AD = DC = \frac{1}{2} AB = 1\)，\(M\) 是 \(PB\) 的中点.

\begin{enumerate}
\item 证明：面 \(PAD \perp\) 面 \(PCD\)；
\item 求 \(AC\) 与 \(PB\) 所成的角；
\item 求面 \(AMC\) 与面 \(BMC\) 所成二面角的大小.
\end{enumerate}


\bitmapfigure[max width=0.42\linewidth]{img/2005/national_paper_1_science/q18_fig1.png}
\end{problem}

\begin{solution}
\begin{enumerate}
\item 由 \(PA\perp\) 底面 \(ABCD\)，知 \(PA\perp CD\). 又因为 \(AB\parallel DC\) 且 \(\angle DAB=90^{\circ}\)，所以 \(AD\perp CD\). 直线 \(PA\)，\(AD\) 是面 \(PAD\) 内相交的两条直线，且 \(CD\perp PA\)，\(CD\perp AD\)，故 \(CD\perp\) 面 \(PAD\). 由于 \(CD\) 在面 \(PCD\) 内，所以面 \(PAD\perp\) 面 \(PCD\).
\item 建立空间直角坐标系，取 \(A\) 为原点，\(AB\)，\(AD\)，\(AP\) 所在方向分别为 \(x\) 轴，\(y\) 轴，\(z\) 轴，则 \(A(0,0,0)\)，\(B(2,0,0)\)，\(D(0,1,0)\)，\(C(1,1,0)\)，\(P(0,0,1)\). 于是
\(\overrightarrow{AC}=(1,1,0)\)，\(\overrightarrow{PB}=(2,0,-1)\).
设 \(AC\) 与 \(PB\) 所成的角为 \(\theta\)，则
\[
\cos\theta=\frac{\left|\overrightarrow{AC}\cdot\overrightarrow{PB}\right|}{\left|\overrightarrow{AC}\right|\left|\overrightarrow{PB}\right|}=\frac{2}{\sqrt{2}\sqrt{5}}=\frac{\sqrt{10}}{5}
\]
所以 \(AC\) 与 \(PB\) 所成的角为 \(\arccos\dfrac{\sqrt{10}}{5}\).
\item 由 \(M\) 是 \(PB\) 的中点，得 \(M(1,0,\frac{1}{2})\). 设 \(N\) 为 \(CM\) 上使 \(AN\perp CM\) 的点. 由
\(\overrightarrow{CM}=(0,-1,\tfrac{1}{2})\)，\(\overrightarrow{CA}=(-1,-1,0)\)，\(\overrightarrow{CB}=(1,-1,0)\)
可得 \(\overrightarrow{CA}\cdot\overrightarrow{CM}=\overrightarrow{CB}\cdot\overrightarrow{CM}=1\). 因此 \(A\)，\(B\) 在直线 \(CM\) 上的射影相同，\(N\) 同时也是 \(B\) 到 \(CM\) 的垂足. 由 \(AN\perp CM\) 得 \(N(1,\frac{1}{5},\frac{2}{5})\)，从而
\(\overrightarrow{NA}=(-1,-\tfrac{1}{5},-\tfrac{2}{5})\)，\(\overrightarrow{NB}=(1,-\tfrac{1}{5},-\tfrac{2}{5})\).
又 \(BN\perp CM\)，故 \(\angle ANB\) 是所求二面角的平面角. 因此
\[
\cos\angle ANB=\frac{\overrightarrow{NA}\cdot\overrightarrow{NB}}{\left|\overrightarrow{NA}\right|\left|\overrightarrow{NB}\right|}=\frac{-1+\frac{1}{25}+\frac{4}{25}}{1+\frac{1}{25}+\frac{4}{25}}=-\frac{2}{3}
\]
所以所求二面角的大小为 \(\arccos\left(-\dfrac{2}{3}\right)\).
\end{enumerate}
\end{solution}

\begin{problem}
设等比数列 \(\{a_{n}\}\) 的公比为 \(q\)，前 \(n\) 项和 \(S_{n} > 0\)（\(n = 1\)，\(2\)，\(\cdots\)）.

\begin{enumerate}
\item 求 \(q\) 的取值范围；
\item 设 \(b_{n} = a_{n+2} - \frac{3}{2} a_{n+1}\)，记 \(\{b_{n}\}\) 的前 \(n\) 项和为 \(T_{n}\)，试比较 \(S_{n}\) 和 \(T_{n}\) 的大小.
\end{enumerate}
\end{problem}

\begin{solution}
\begin{enumerate}
\item 因为 \(S_1=a_1>0\)，所以 \(a_1>0\). 当 \(q\ne1\) 时，
\[
S_n=a_1\frac{1-q^n}{1-q}
\]
若 \(q\geqslant0\)，当 \(0\leqslant q<1\) 时分子、分母均为正，当 \(q>1\) 时分子、分母均为负，且 \(q=1\) 时 \(S_n=na_1>0\)，所以所有 \(q\geqslant0\) 均满足条件. 若 \(q<0\)，则奇数 \(n\) 时 \(1-q^n>0\)，而偶数 \(n\) 时必须有 \(1-q^n>0\)，即 \(|q|<1\). 因此 \(q\) 的取值范围为 \(q>-1\).
\item 由 \(a_{n+1}=a_1q^n\)，得
\[
b_n=a_1q^n\left(q-\frac{3}{2}\right)
\]
于是
\[
T_n=\left(q-\frac{3}{2}\right)\sum_{k=1}^{n}a_1q^k=q\left(q-\frac{3}{2}\right)S_n
\]
因此
\[
S_n-T_n=\left[1-q\left(q-\frac{3}{2}\right)\right]S_n=\left(2-q\right)\left(q+\frac{1}{2}\right)S_n
\]
由于 \(S_n>0\)，所以
\[
\begin{cases}
S_n<T_n,&-1<q<-\dfrac{1}{2}\text{ 或 }q>2,\\
S_n=T_n,&q=-\dfrac{1}{2}\text{ 或 }q=2,\\
S_n>T_n,&-\dfrac{1}{2}<q<2.
\end{cases}
\]
\end{enumerate}
\end{solution}

\begin{problem}
\(9\) 粒种子分种在 \(3\) 个坑内，每坑 \(3\) 粒. 每粒种子发芽的概率为 \(0.5\)，若一个坑内至少有 \(1\) 粒种子发芽，则这个坑不需要补种，若一个坑里的种子都没发芽，则这个坑需要补种，假定每个坑至多补种一次. 每补种 \(1\) 个坑需 \(10\text{ 元}\)，用 \(\xi\) 表示补种费用. 写出 \(\xi\) 的分布列并求 \(\xi\) 的数学期望.（精确到 \(0.01\)）
\end{problem}

\begin{solution}
一个坑内的 \(3\) 粒种子都不发芽的概率为 \(\left(1-0.5\right)^3=\dfrac{1}{8}\)，一个坑需要补种的概率为 \(\dfrac{1}{8}\)，不需要补种的概率为 \(\dfrac{7}{8}\). 设需要补种的坑数为 \(Y\)，则 \(Y\) 服从参数为 \(3\)，\(\dfrac{1}{8}\) 的二项分布，且 \(\xi=10Y\). 因此
\(P(\xi=10k)=\mathrm C_3^k\left(\frac{1}{8}\right)^k\left(\frac{7}{8}\right)^{3-k}\)，\((k=0,1,2,3)\).
故 \(\xi\) 的分布列为
\begin{center}
\begin{tabular}{c|cccc}
\(\xi\) & \(0\) & \(10\) & \(20\) & \(30\) \\
\hline
\(P\) & \(\dfrac{343}{512}\) & \(\dfrac{147}{512}\) & \(\dfrac{21}{512}\) & \(\dfrac{1}{512}\)
\end{tabular}
\end{center}
于是
\[
E(\xi)=10E(Y)=10\cdot3\cdot\frac{1}{8}=\frac{15}{4}=3.75\text{ 元}
\]
所以 \(\xi\) 的数学期望为 \(3.75\text{ 元}\).
\end{solution}

\begin{problem}
已知椭圆的中心为坐标原点 \(O\)，焦点在 \(x\) 轴上，斜率为 \(1\) 且过椭圆右焦点 \(F\) 的直线交椭圆于 \(A\)，\(B\) 两点，\(\overrightarrow{OA} + \overrightarrow{OB}\) 与 \(\bs{a} = (3, -1)\) 共线.

\begin{enumerate}
\item 求椭圆的离心率；
\item 设 \(M\) 为椭圆上任意一点，且 \(\overrightarrow{OM} = \lambda \overrightarrow{OA} + \mu \overrightarrow{OB}\)（\(\lambda\)，\(\mu \in \R\)），证明 \(\lambda^2 + \mu^2\) 为定值.
\end{enumerate}
\end{problem}

\begin{solution}
\begin{enumerate}
\item 设椭圆方程为 \(\dfrac{x^2}{a_0^2}+\dfrac{y^2}{b_0^2}=1\)，其中 \(a_0>b_0>0\)，焦半距为 \(c\)，则 \(c^2=a_0^2-b_0^2\)，右焦点为 \(F(c,0)\). 由直线斜率为 \(1\)，该直线方程为 \(y=x-c\). 设它与椭圆的交点为 \(A(x_1,y_1)\)，\(B(x_2,y_2)\)，代入椭圆方程得
\[
(a_0^2+b_0^2)x^2-2a_0^2cx+a_0^2(c^2-b_0^2)=0
\]
由韦达定理，
\(x_1+x_2=\frac{2a_0^2c}{a_0^2+b_0^2}\)，\(x_1x_2=\frac{a_0^2(c^2-b_0^2)}{a_0^2+b_0^2}\).
又 \(y_i=x_i-c\)，所以
\[
y_1+y_2=-\frac{2b_0^2c}{a_0^2+b_0^2}
\]
因 \(\overrightarrow{OA}+\overrightarrow{OB}\) 与 \((3,-1)\) 共线，故
\[
\frac{y_1+y_2}{x_1+x_2}=-\frac{1}{3}
\]
因此 \(-\dfrac{b_0^2}{a_0^2}=-\dfrac{1}{3}\)，即 \(b_0^2=\dfrac{a_0^2}{3}\). 于是
\[
e=\frac{c}{a_0}=\sqrt{1-\frac{b_0^2}{a_0^2}}=\sqrt{\frac{2}{3}}=\frac{\sqrt{6}}{3}
\]
\item 由 \(b_0^2=\dfrac{a_0^2}{3}\) 和 \(c^2=\dfrac{2a_0^2}{3}\)，可得
\(x_1+x_2=\frac{3c}{2}\)，\(x_1x_2=\frac{a_0^2}{4}\).
所以
\[
(x_1-x_2)^2=(x_1+x_2)^2-4x_1x_2=\frac{a_0^2}{2}
\]
交换 \(A\)，\(B\) 只会交换 \(\lambda\)，\(\mu\)，不影响 \(\lambda^2+\mu^2\)，故不妨设 \(x_1-x_2=\dfrac{a_0}{\sqrt{2}}\). 因为 \(y_i=x_i-c\)，于是
\(\overrightarrow{OA}+\overrightarrow{OB}=\left(\frac{3c}{2},-\frac{c}{2}\right)\)，\(\overrightarrow{OA}-\overrightarrow{OB}=\left(\frac{a_0}{\sqrt{2}},\frac{a_0}{\sqrt{2}}\right)\).
记 \(s=\lambda+\mu\)，\(d=\lambda-\mu\)，则
\[
\overrightarrow{OM}=\frac{s}{2}\left(\overrightarrow{OA}+\overrightarrow{OB}\right)+\frac{d}{2}\left(\overrightarrow{OA}-\overrightarrow{OB}\right)
\]
因此 \(M\) 的坐标为
\(x_M=\frac{3cs}{4}+\frac{a_0d}{2\sqrt{2}}\)，\(y_M=-\frac{cs}{4}+\frac{a_0d}{2\sqrt{2}}\).
令 \(u=\dfrac{cs}{4}\)，\(v=\dfrac{a_0d}{2\sqrt{2}}\)，则由 \(M\) 在椭圆上以及 \(b_0^2=\dfrac{a_0^2}{3}\)，有
\[
1=\frac{x_M^2}{a_0^2}+\frac{y_M^2}{b_0^2}=\frac{(3u+v)^2+3(-u+v)^2}{a_0^2}=\frac{12u^2+4v^2}{a_0^2}=\frac{s^2+d^2}{2}
\]
而 \(\dfrac{s^2+d^2}{2}=\dfrac{(\lambda+\mu)^2+(\lambda-\mu)^2}{2}=\lambda^2+\mu^2\)，所以 \(\lambda^2+\mu^2=1\)，为定值.
\end{enumerate}
\end{solution}

\begin{problem}
\begin{enumerate}
\item 设函数 \(f(x) = x\log_2 x + (1 - x)\log_2(1 - x)\)（\(0 < x < 1\)），求 \(f(x)\) 的最小值；
\item 设正数 \(p_1\)，\(p_2\)，\(p_3\)，\(\cdots\)，\(p_{2^n}\) 满足 \(p_1 + p_2 + p_3 + \cdots + p_{2^n} = 1\)，证明 \(p_1 \log_2 p_1 + p_2 \log_2 p_2 + p_3 \log_2 p_3 + \cdots + p_{2^n} \log_2 p_{2^n} \geqslant -n\).
\end{enumerate}
\end{problem}

\begin{solution}
\begin{enumerate}
\item 在 \(0<x<1\) 内，
\[
f'(x)=\log_2x-\log_2(1-x)=\log_2\frac{x}{1-x}
\]
所以 \(f'(x)=0\) 当且仅当 \(x=\dfrac{1}{2}\). 又
\[
f''(x)=\frac{1}{\ln2}\left(\frac{1}{x}+\frac{1}{1-x}\right)>0
\]
故 \(f(x)\) 在 \(x=\dfrac{1}{2}\) 处取得最小值. 最小值为
\[
f\left(\frac{1}{2}\right)=\frac{1}{2}\log_2\frac{1}{2}+\frac{1}{2}\log_2\frac{1}{2}=-1
\]
\item 用数学归纳法证明. \(n=1\) 时，令 \(p_1=x\)，\(p_2=1-x\)，由第（1）问得
\[
p_1\log_2p_1+p_2\log_2p_2\geqslant-1
\]
假设对 \(2^{n-1}\) 个正数 \(q_1\)，\(q_2\)，\(\cdots\)，\(q_{2^{n-1}}\)，当 \(\sum_{j=1}^{2^{n-1}}q_j=1\) 时，不等式
\[
\sum_{j=1}^{2^{n-1}}q_j\log_2q_j\geqslant-(n-1)
\]
成立. 对 \(2^n\) 个正数作分组，令
\(q_j=p_{2j-1}+p_{2j}\)，\(t_j=\frac{p_{2j-1}}{q_j}\)，\(\left(j=1,2,\cdots,2^{n-1}\right)\).
则 \(0<t_j<1\)，\(\sum q_j=1\)，并且
\[
p_{2j-1}\log_2p_{2j-1}+p_{2j}\log_2p_{2j}=q_j\log_2q_j+q_j\left[t_j\log_2t_j+(1-t_j)\log_2(1-t_j)\right]\geqslant q_j\log_2q_j-q_j
\]
对 \(j=1\)，\(2\)，\(\cdots\)，\(2^{n-1}\) 求和，得到
\[
\sum_{i=1}^{2^n}p_i\log_2p_i\geqslant\sum_{j=1}^{2^{n-1}}q_j\log_2q_j-\sum_{j=1}^{2^{n-1}}q_j\geqslant-(n-1)-1=-n
\]
因此不等式对任意正整数 \(n\) 成立.
\end{enumerate}
\end{solution}
